\documentclass[11pt,letterpaper]{article}
% Reusable style. Build each .tex from its own directory using scripts/build.sh.
\usepackage[T1]{fontenc}
\usepackage[utf8]{inputenc}
\usepackage{lmodern}
\usepackage[margin=0.86in,headheight=15pt]{geometry}
\usepackage{amsmath,amssymb,amsthm,mathtools}
\usepackage{booktabs,longtable,tabularx,array,multicol}
\usepackage[table]{xcolor}
\usepackage{microtype}
\usepackage{enumitem}
\usepackage{fancyhdr}
\usepackage{titlesec}
\usepackage{hyperref}
\usepackage{listings}
\usepackage[most]{tcolorbox}
\usepackage{needspace}
\definecolor{aeNavy}{HTML}{172D43}
\definecolor{aeBlue}{HTML}{176B90}
\definecolor{aeTeal}{HTML}{148477}
\definecolor{aePale}{HTML}{F2F7F9}
\definecolor{aeGray}{HTML}{5C6771}
\hypersetup{colorlinks=true,linkcolor=aeBlue,urlcolor=aeBlue,pdfauthor={Algorithm Engineering Documentation Suite},pdfsubject={Independent educational materials based on the supplied TAOCP volume topics}}
\titleformat{\section}{\Large\bfseries\color{aeNavy}}{\thesection}{0.6em}{}
\titleformat{\subsection}{\large\bfseries\color{aeBlue}}{\thesubsection}{0.6em}{}
\titleformat{\subsubsection}{\normalsize\bfseries\color{aeNavy}}{\thesubsubsection}{0.6em}{}
\setlist[itemize]{leftmargin=*,topsep=3pt,itemsep=1pt}
\setlist[enumerate]{leftmargin=*,topsep=3pt,itemsep=3pt}
\setlength{\parindent}{0pt}\setlength{\parskip}{5pt}
\pagestyle{fancy}\fancyhf{}\fancyhead[L]{\footnotesize\color{aeGray}ALGORITHM ENGINEERING}\fancyhead[R]{\footnotesize\color{aeGray}DOCUMENTATION SUITE}\fancyfoot[L]{\footnotesize\color{aeGray}Independent study and engineering practice}\fancyfoot[R]{\footnotesize\thepage}
\renewcommand{\headrulewidth}{0.35pt}
\newtcolorbox{engineeringnote}[1][]{breakable,colback=aePale,colframe=aeBlue,boxrule=0.5pt,arc=1mm,fonttitle=\bfseries,title={Engineering note},#1}
\newtcolorbox{evidencebox}[1][]{breakable,colback=white,colframe=aeTeal,boxrule=0.6pt,arc=1mm,fonttitle=\bfseries,title={Evidence and deliverables},#1}
\lstdefinestyle{algo}{basicstyle=\ttfamily\footnotesize,breaklines=true,frame=single,rulecolor=\color{aeGray!40},backgroundcolor=\color{aePale},showstringspaces=false,tabsize=2,columns=fullflexible,keepspaces=true,keywordstyle=\color{aeBlue}\bfseries,commentstyle=\color{aeTeal}}
\lstset{style=algo}
\newcommand{\DocTitle}[3]{%
\hypersetup{pageanchor=false}\begin{titlepage}\vspace*{1.2cm}
{\color{aeTeal}\rule{\linewidth}{2.3pt}}\par\vspace{1.15cm}
{\fontsize{24}{30}\selectfont\bfseries\color{aeNavy}\raggedright #1\par}\vspace{0.45cm}
{\Large\color{aeBlue}#2\par}\vspace{1.3cm}
{\large\bfseries #3\par}\vspace{0.35cm}
{\color{aeGray}Edition 1.0 \quad | \quad 29 October 2026\par}\vfill
{\color{aeGray}\small Original educational material aligned to the chapter and section topics in the five supplied volumes of Donald E. Knuth's \emph{The Art of Computer Programming}. This document is not an authorized reproduction or substitute for the original books.}\par\vspace{0.4cm}
{\color{aeTeal}\rule{\linewidth}{0.75pt}}
\end{titlepage}\hypersetup{pageanchor=true}\pagenumbering{roman}}
\newcommand{\SourceNote}[1]{\par\smallskip{\footnotesize\color{aeGray}\textbf{Source locator:} #1}\par\smallskip}
\newcommand{\Outcome}[1]{\par\textbf{Competency target.} #1\par}
\newcommand{\Checkpoint}[1]{\par\textbf{Checkpoint.} #1\par}
% Intentionally simple printable writing lines: call \AnswerLines{number}.
\newcount\alcount
\newcommand{\AnswerLines}[1]{\par\alcount=0\loop\ifnum\alcount<#1\par\noindent\rule{\linewidth}{0.3pt}\vspace{0.33cm}\advance\alcount by1\repeat\par}
\newcommand{\FormField}[2][0.62\linewidth]{\noindent\textbf{#2}\hspace{0.65em}\rule{#1}{0.4pt}\par\vspace{0.18cm}}
\newcommand{\ReviewTable}{\begin{tabularx}{\linewidth}{@{}p{0.23\linewidth}X p{0.2\linewidth}@{}}\toprule Criterion & Recorded evidence & Rating (0--4)\\\midrule Correctness & & \\[1em] Proof quality & & \\[1em] Complexity & & \\[1em] Testing & & \\[1em] Reproducibility & & \\[1em] Communication & & \\[1em]\bottomrule\end{tabularx}}

\begin{document}
\DocTitle{Volume 1 Study Guide}{Chapters 1--2: foundations and information structures}{Chapter- and section-aligned study guide}
\tableofcontents
\clearpage
\pagenumbering{arabic}
\section{Study protocol}

For each subsection, read the specified part of the original volume, keep a mathematical notebook, produce an original implementation or design sketch, and evaluate the acceptance gate. The investigations below are new exercises, not copies of the book's problems or solutions. Annotate any starred or mathematically advanced source sections as optional on the first pass; revisit them after the prerequisite mathematics is in place.
\begin{evidencebox}Recommended artifacts: study notes, problem specification, derivation, runnable code or model, tests, performance evidence and reflective comparison to original assumptions.\end{evidencebox}

\section{\S 1.1: Algorithms}
\SourceNote{Knuth, Volume 1, \S 1.1.}
\Outcome{Define deterministic process, input/output contracts and finite procedure. Distinguish a mathematical relation from an executable method.}
\subsection{Independent investigation}
Specify an iterative greatest-common-divisor routine for nonnegative integers; include zero cases.
\subsection{Analysis and verification gate}
Prove its loop measure strictly decreases and write assertions for both return and no-mutation properties.
\subsection{Study questions}
\begin{enumerate}
\item Which representations, input assumptions and cost metrics control this technique's applicability?
\item Which invariant or combinatorial identity establishes correctness, and which input is its strongest counterexample candidate?
\item What result would cause the implementation to fail review despite passing a small collection of examples?
\end{enumerate}
\Checkpoint{Preserve a short explanation, the relevant test evidence, and one unresolved question before proceeding.}

\section{\S 1.2.1--1.2.4: Induction, sums, number theory}
\SourceNote{Knuth, Volume 1, \S 1.2.1--1.2.4.}
\Outcome{Use strong induction, logarithmic transformations, sum identities and integer division to reason about cost.}
\subsection{Independent investigation}
Analyze nested loops with triangular iteration counts and test the exact iteration formula.
\subsection{Analysis and verification gate}
Prove upper and lower quadratic bounds with explicit constants and starting input size.
\subsection{Study questions}
\begin{enumerate}
\item Which representations, input assumptions and cost metrics control this technique's applicability?
\item Which invariant or combinatorial identity establishes correctness, and which input is its strongest counterexample candidate?
\item What result would cause the implementation to fail review despite passing a small collection of examples?
\end{enumerate}
\Checkpoint{Preserve a short explanation, the relevant test evidence, and one unresolved question before proceeding.}

\section{\S 1.2.5--1.2.9: Permutations, binomial coefficients, recurrences}
\SourceNote{Knuth, Volume 1, \S 1.2.5--1.2.9.}
\Outcome{Relate factorial enumeration, combination counts, harmonic sums, Fibonacci recurrences and generating functions.}
\subsection{Independent investigation}
Derive the number of leaves of a binary inclusion/exclusion tree and write a lazy enumerator.
\subsection{Analysis and verification gate}
Explain why output cardinality imposes a lower bound independent of constant-factor optimization.
\subsection{Study questions}
\begin{enumerate}
\item Which representations, input assumptions and cost metrics control this technique's applicability?
\item Which invariant or combinatorial identity establishes correctness, and which input is its strongest counterexample candidate?
\item What result would cause the implementation to fail review despite passing a small collection of examples?
\end{enumerate}
\Checkpoint{Preserve a short explanation, the relevant test evidence, and one unresolved question before proceeding.}

\section{\S 1.2.10--1.2.11: Algorithm analysis and asymptotic expansions}
\SourceNote{Knuth, Volume 1, \S 1.2.10--1.2.11.}
\Outcome{Keep exact operation counts separate from asymptotic notation and measured runtime.}
\subsection{Independent investigation}
Count comparisons in a linear scan for successful and unsuccessful queries on length n inputs.
\subsection{Analysis and verification gate}
Provide best, worst and distribution-conditioned expected counts with justified assumptions.
\subsection{Study questions}
\begin{enumerate}
\item Which representations, input assumptions and cost metrics control this technique's applicability?
\item Which invariant or combinatorial identity establishes correctness, and which input is its strongest counterexample candidate?
\item What result would cause the implementation to fail review despite passing a small collection of examples?
\end{enumerate}
\Checkpoint{Preserve a short explanation, the relevant test evidence, and one unresolved question before proceeding.}

\section{\S 1.3: MIX and instruction-level models}
\SourceNote{Knuth, Volume 1, \S 1.3.}
\Outcome{Recognize MIX as the historic model used by the source for expressing machine-oriented algorithms.}
\subsection{Independent investigation}
Produce a mapping from an abstract array operation to plausible load, compare and branch operations.
\subsection{Analysis and verification gate}
Identify which MIX assumptions do not transfer directly to modern caches or compiler optimizations.
\subsection{Study questions}
\begin{enumerate}
\item Which representations, input assumptions and cost metrics control this technique's applicability?
\item Which invariant or combinatorial identity establishes correctness, and which input is its strongest counterexample candidate?
\item What result would cause the implementation to fail review despite passing a small collection of examples?
\end{enumerate}
\Checkpoint{Preserve a short explanation, the relevant test evidence, and one unresolved question before proceeding.}

\section{\S 1.4: Subroutines, coroutines, interpreters and I/O}
\SourceNote{Knuth, Volume 1, \S 1.4.}
\Outcome{Compare control flow, calling conventions, transfer of state and deterministic I/O.}
\subsection{Independent investigation}
Write an iterator-style stream processor and explain state persistence between yields.
\subsection{Analysis and verification gate}
Test repeated resume operations, exhausted iteration and failure recovery.
\subsection{Study questions}
\begin{enumerate}
\item Which representations, input assumptions and cost metrics control this technique's applicability?
\item Which invariant or combinatorial identity establishes correctness, and which input is its strongest counterexample candidate?
\item What result would cause the implementation to fail review despite passing a small collection of examples?
\end{enumerate}
\Checkpoint{Preserve a short explanation, the relevant test evidence, and one unresolved question before proceeding.}

\section{\S 2.1--2.2: Linear lists}
\SourceNote{Knuth, Volume 1, \S 2.1--2.2.}
\Outcome{Compare contiguous, singly linked, doubly linked and circular representations and queue operations.}
\subsection{Independent investigation}
Implement a bounded ring buffer; measure enqueue/dequeue under varying occupancy.
\subsection{Analysis and verification gate}
Demonstrate empty/full/wraparound behavior and derive amortized or worst-case bounds.
\subsection{Study questions}
\begin{enumerate}
\item Which representations, input assumptions and cost metrics control this technique's applicability?
\item Which invariant or combinatorial identity establishes correctness, and which input is its strongest counterexample candidate?
\item What result would cause the implementation to fail review despite passing a small collection of examples?
\end{enumerate}
\Checkpoint{Preserve a short explanation, the relevant test evidence, and one unresolved question before proceeding.}

\section{\S 2.3.1--2.3.4: Trees and traversals}
\SourceNote{Knuth, Volume 1, \S 2.3.1--2.3.4.}
\Outcome{Distinguish binary tree traversal orders, ordered-tree encoding, structural properties and path length.}
\subsection{Independent investigation}
Generate small trees and verify preorder/inorder/postorder traversals using a reference traversal.
\subsection{Analysis and verification gate}
Prove traversal visits every node once without revisiting any node.
\subsection{Study questions}
\begin{enumerate}
\item Which representations, input assumptions and cost metrics control this technique's applicability?
\item Which invariant or combinatorial identity establishes correctness, and which input is its strongest counterexample candidate?
\item What result would cause the implementation to fail review despite passing a small collection of examples?
\end{enumerate}
\Checkpoint{Preserve a short explanation, the relevant test evidence, and one unresolved question before proceeding.}

\section{\S 2.3.5--2.5: Memory management and multilinked structures}
\SourceNote{Knuth, Volume 1, \S 2.3.5--2.5.}
\Outcome{Specify ownership, cross-links, lifetime, freelists and allocation costs.}
\subsection{Independent investigation}
Compare a vector-backed adjacency representation against pointer-linked nodes under the same workload.
\subsection{Analysis and verification gate}
Document memory overhead, traversal locality and iterator invalidation.
\subsection{Study questions}
\begin{enumerate}
\item Which representations, input assumptions and cost metrics control this technique's applicability?
\item Which invariant or combinatorial identity establishes correctness, and which input is its strongest counterexample candidate?
\item What result would cause the implementation to fail review despite passing a small collection of examples?
\end{enumerate}
\Checkpoint{Preserve a short explanation, the relevant test evidence, and one unresolved question before proceeding.}

\section{Volume synthesis and mastery review}
\begin{enumerate}
\item Draw a dependency graph connecting at least four major ideas from this volume.
\item Select two alternative implementations for one workload and compare their preconditions, proven cost, memory use, measured performance and failure cases.
\item Present a five-minute technical defense of the most surprising experimental result. A claim that differs from the mathematical model requires a measured or documented explanation.
\end{enumerate}
\subsection{Completion checklist}
\begin{itemize}\item All module notes include a source locator and an independently stated learning outcome.
\item Every implementation has reproducible tests and edge cases.
\item Every complexity claim identifies its cost model and underlying assumptions.
\item A peer can independently reproduce one chosen experiment using its recorded configuration.
\end{itemize}
\end{document}
